% from ega3-1-fr.tex
\oldpage[0\textsubscript{III}]{19}
\begin{env}[10.2.3]
\label{0.10.2.3-fr}
Les hypothèses sur $A$, $B$, $\mathfrak{J}$, $M$ étant celles formulées au
début de (\hyperref[0.10.2.2-fr]{10.2.2}), soient $\widehat A$ le séparé
complété de $A$ pour la topologie $\mathfrak{J}$-pré-adique, $\widehat M$ le
séparé complété de $M$ pour la topologie $\mathfrak{J}B$-pré-adique. Alors,
pour que $M$ soit un $A$-module plat, il faut et il suffit que $\widehat M$
soit un $\widehat A$-module plat.
\end{env}

\begin{env}[10.2.4]
\label{0.10.2.4-fr}
Soient $\rho:A\to B$ un homomorphisme local d'anneaux locaux noethériens, $k$
le corps résiduel de $A$, $M$, $N$ deux $B$-modules de type fini, $N$ étant
supposé être $A$-\emph{plat}. Soit $u:M\to N$ un $B$-homomorphisme. Alors les
conditions suivantes sont équivalentes~:
\begin{enumerate}
  \item[$a)$] $u$ est injectif et $\operatorname{Coker}(u)$ est un $A$-module plat.
  \item[$b)$] $u\otimes1:M\otimes_A k\to N\otimes_A k$ est injectif.
\end{enumerate}
\end{env}

\begin{env}[10.2.5]
\label{0.10.2.5-fr}
Soient $\rho:A\to B$, $\sigma:B\to C$ des homomorphismes locaux d'anneaux
locaux noethériens, $k$ le corps résiduel de $A$, $M$ un $C$-module de type
fini. On suppose que $B$ est un $A$-module \emph{plat} Alors les conditions
suivantes sont équivalentes~:
\begin{enumerate}
  \item[$a)$] $M$ est un $B$-module plat.
  \item[$b)$] $M$ est un $A$-module plat, et $M\otimes_A k$ est un
    $(B\otimes_A k)$-module plat.
\end{enumerate}
\end{env}

\begin{proposition}[10.2.6]
\label{0.10.2.6-fr}
Soient $A$, $B$ deux anneaux locaux noethériens, $\rho:A\to B$ un
homomorphisme local, $\mathfrak{J}$ un idéal de $B$ contenu dans l'idéal
maximal, $M$ un $B$-module de type fini. Supposons que pour tout $n\geq0$,
$M_n=M/\mathfrak{J}^{n+1}M$ soit un $A$-module plat. Alors $M$ est un
$A$-module plat.
\end{proposition}

Il faut prouver que pour tout homomorphisme injectif $u:N'\to N$ de
$A$-modules de type fini, $v=1\otimes u:M\otimes_A N'\to M\otimes_A N$ est
injectif. Or, $M\otimes_A N'$ et $M\otimes_A N$ sont des $B$-modules de type
fini, donc séparés pour la topologie $\mathfrak{J}$-pré-adique
(\hyperref[0.7.3.5-fr]{$0_{\mathrm I}$, 7.3.5})~; il suffit donc de prouver que
l'homomorphisme
$\widehat v:\widehat{(M\otimes_A N')}\to\widehat{(M\otimes_A N)}$ pour les
séparés complétés est injectif. Or, on a $\widehat v=\varprojlim v_n$, où
$v_n$ est l'homomorphisme $1\otimes u:M_n\otimes_A N'\to M_n\otimes_A N$~;
comme par hypothèse $M_n$ est $A$-plat, $v_n$ est injectif pour tout $n$, donc
il en est de même de $v$, le foncteur $\varprojlim$ étant exact à gauche.

\begin{corollary}[10.2.7]
\label{0.10.2.7-fr}
Soient $A$ un anneau noethérien, $B$ un anneau local noethérien,
$\rho:A\to B$ un homomorphisme, $f$ un élément de l'idéal maximal de $B$, $M$
un $B$-module de type fini. Supposons que l'homothétie
$f_M:x\mapsto fx$ de $M$ soit injective et que $M/fM$ soit un $A$-module
plat. Alors $M$ est un $A$-module plat.
\end{corollary}

Posons $M_i=f^iM$ pour $i\geq0$~; comme $f_M$ est injective,
$M_i/M_{i+1}$ est isomorphe à $M/fM$, donc $A$-plat pour tout $i\geq0$~; de
la suite exacte
\[
  0\longrightarrow M_i/M_{i+1}\longrightarrow M/M_{i+1}
  \longrightarrow M/M_i\longrightarrow0
\]
on tire par récurrence sur $i$ que $M/M_i$ est $A$-\emph{plat} pour tout
$i\geq0$ (\hyperref[0.6.1.2-fr]{$0_{\mathrm I}$, 6.1.2})~; on peut donc
appliquer (\hyperref[0.10.2.6-fr]{10.2.6}). On peut aussi raisonner directement
comme suit~: pour tout $A$-module $N$ de type fini, $M\otimes_A N$ est un
$B$-module de type fini~; comme $f$ appartient au radical $\mathfrak n$ de
$B$, la topologie $(f)$-adique sur $M\otimes_A N$ est plus fine que la
